


\begin{figure}[t]
    \centering
    \includegraphics[width=0.7\linewidth]{images/web-tool-table.png}
    \caption{InferQ's web-based UI for exploring dataset and circuit features. }
    \label{fig:inferq-web}
     % \vspace{-0.2cm}
\end{figure}


 

\section{Implementation and Artifacts}
\label{sec:impl}

\sys ships as (i) a Python framework that turns a user configuration
(Table~\ref{tab:base_params}) into \emph{circuits, features, and an RDBMS-ready SQL
workload} for simulation, and (ii) a web-based UI for interactive dataset exploration
(Fig.~\ref{fig:inferq-web}). Concretely, a database researcher can provide a small
configuration file (e.g., qubit count) and run the toolkit to obtain  SQL
queries (plus circuit artifacts in JSON and QPY formats) that are ready to execute in an
RDBMS engine. \revisionM{Detailed descriptions of all circuit templates and the corresponding SQL code are provided in the online documentation.\footnote{\url{https://github.com/InfiniData-Lab/InferQ/blob/main/generators/sql-templates.md}}}
In addition to on-demand generation, \sys releases a large dataset of \num{202975} circuits online.\footnote{\url{https://github.com/InfiniData-Lab/InferQ/blob/main/analysis/training_data/estimator_training_data.parquet}}
Both database researchers and quantum scientists can search, filter, and download
circuits and feature records through the web UI.

The framework implements the pipeline described in
Section~\ref{sec:benchmark} and Section~\ref{ssec:features}.
% It (1) generates a circuit by instantiating a sequence of circuit templates, (2) emits the corresponding SQL workload for
% simulation, and (3) extracts circuit features.
Each circuit is assigned a \emph{content-based hash} computed from its structure and
sampled parameters. This hash serves as the primary key and links the circuit artifact,
the SQL workload, and all extracted feature records. For each circuit instance, \sys stores:
(i) the generated SQL queries (as \texttt{.sql} scripts) that represent
the simulation workload,
(ii) the circuit in a portable Qiskit serialization format (e.g., \texttt{.qpy}), and
(iii) a JSON record containing the template provenance and extracted features. 
\revisionB{InferQ includes a metadata analysis function that aggregates circuit statistics from JSON files, such as gate mix, depth, width, and interaction-graph properties. Details are provided in Appendix~J of \cite{inferqtech} and in the online repository.\footnote{\url{https://github.com/InfiniData-Lab/InferQ/blob/main/analysis/distributions/distributions.py}}}
% \ref{app:circuit-characterization}

 \medskip
\para{Reproducibility} The code of our benchmark can be found online.\footnote{\url{https://github.com/InfiniData-Lab/InferQ/blob/main/README.md}}
\sys is designed to be reproducible. Given the same code revision, dependency
lockfile, and global seed (Table~\ref{tab:base_params}), the framework deterministically
reproduces the same template sequence, sampled parameters, circuit artifact, emitted SQL,
and extracted features. The content-based hash and provenance record make each benchmark
instance easy to reference, regenerate, and audit. We include more implementation and reproducibility
details in our technical report \cite{inferqtech}.



